% !TeX spellcheck = it_IT
\title{Schema ComplAlgo}
\author{Massimo Perego}
\date{}

\documentclass[11pt]{article}
\usepackage{graphicx} 
\usepackage{amsmath}
\usepackage{amssymb}
\usepackage{amsfonts}
\usepackage[hidelinks]{hyperref}
\usepackage[autostyle, english = american]{csquotes}
\usepackage[parfill]{parskip}
\MakeOuterQuote{"}

\usepackage{bm}

\renewcommand{\contentsname}{Indice}

\renewcommand{\O}{\mathcal{O}}
\renewcommand{\P}{\mathcal{P}}
\renewcommand{\Pr}{\mathbb{P}}
\renewcommand{\H}{\mathcal{H}}

\newcommand{\NP}{\mathcal{NP}}
\newcommand{\coNP}{\text{co-}\NP}
\newcommand{\coP}{\text{co-}\P}
\newcommand{\NPc}{$\NP$-completo}
\newcommand{\BPP}{\mathcal{BPP}}
\newcommand{\coBPP}{\text{co-}\BPP}
\newcommand{\RP}{\mathcal{RP}}
\newcommand{\coRP}{\text{co-}\RP}
\newcommand{\ZPP}{\mathcal{ZPP}}
\newcommand{\I}{\mathcal{I}}
\newcommand{\X}{\mathcal{X}}
\newcommand{\Y}{\mathcal{Y}}
\newcommand{\C}{\mathcal{C}}
\newcommand{\A}{\mathcal{A}}
\newcommand{\U}{\mathcal{U}}
\newcommand{\F}{\mathcal{F}}
\newcommand{\T}{\mathcal{T}}
\newcommand{\E}{\mathcal{E}}
\newcommand{\Nc}{\mathcal{N}}
\newcommand{\Ind}{\mathbb{I}}
\newcommand{\Ex}{\mathbb{E}}
\newcommand{\N}{\mathbb{N}}
\newcommand{\R}{\mathbb{R}}
\newcommand{\polred}{\preceq_\text{p}}
\newcommand{\equivpol}{\equiv_\text{p}}

\newcommand{\Var}{\textnormal{Var}}
\newcommand{\Opt}{\textnormal{OPT}}

\newcommand{\Count}[1]{\textnormal{\textsc{Count}}(#1)}
\newcommand{\Inc}[1]{\textnormal{\textsc{Inc}}(#1)}
\newcommand{\SelEl}[1]{\textnormal{\textsc{SelectEl}}(#1)}
\newcommand{\CMS}[2]{\textnormal{\textsc{CMS}}[#1][#2]}

\begin{document}
	\maketitle
	\tableofcontents
	\newpage	
	
	\section{Complessità Computazionale}
	% !TeX spellcheck = it_IT
% !TeX root = ../scompl.tex

\subsection{Riducibilità polinomiale tra problemi di decisione}

Definizioni di:
\begin{itemize}
	\item Problema di decisione: istanze, funzione di decisione, lunghezza della codifica

	\item Tempo di esecuzione di un problema, oracolo

	\item Riduzione polinomiale
\end{itemize}

Esempi di riducibilità tra problemi:
\begin{itemize}
	\item Independent Set $\equivpol$ Vertex Cover: se $S$ insieme indipendente, il complemento è una copertura e viceversa

	\item Vertex Cover $\polred$ Set Cover: costruire un'istanza di Set Cover con $\U \equiv E$ e $\mathcal{S} \equiv \left\{S_i \mid i \in V \right\}$

	\item Independent Set $\equivpol$ Set Packing: $V \equiv \left\{v_S \mid S \in \mathcal{S} \right\}$, $(v_S, v_T) \in E \Leftrightarrow S \cap T \not \equiv \emptyset$

	\item 3-SAT $\polred$ Independent Set: da formula, grafo con clique di 3 e arco tra ogni letterale e la sua negazione
\end{itemize}

\subsection{Classi $\P$ e $\NP$}

Definizioni di:
\begin{itemize}
	\item Classe $\P$: problemi risolvibili in tempo polinomiale

	\item Classe $\NP$: problemi che posseggono un certificatore polinomiale
\end{itemize}

Rapporto tra $\P$ e $\NP$: $\P \subseteq \NP$, $\P \stackrel{?}{\equiv} \NP$.

\subsection{Problemi $\NP$-completi}

Definizione di \NPc\ e dimostrare $x \in \P \Leftrightarrow \P \equiv \NP$.

Teorema di Cook-Levin:
\begin{itemize}
	\item Definizione di circuito e Circuit Satisfiability

	\item CS è \NPc

	\item 2-IS $\polred$ CS: $\binom{n}{2} + n$ nodi input, due sottoalberi, uno per controllare che sia un IS, uno per controllare ci siano almeno due elementi, messi assieme alla radice,
\end{itemize}

Dimostrazioni:
\begin{itemize}
	\item $Y$ \NPc\ e $X \in \NP$ t.c. $Y \polred X \implies X$ \NPc

	\item CS $\polred$ 3-SAT: formula booleana, circuito, CNF, 3-CNF

	\item \textit{[Opzionale]} 2-SAT $\in \P$
\end{itemize}

Definizioni di
\begin{itemize}
	\item $\coNP$: quando non si sa se $\bar q(I) = 1$
	\begin{itemize}
		\item $\coNP \not \equiv \NP \implies \P \not \equiv \NP$
	\end{itemize}

	\item $\coP$: invertire $q$
	\begin{itemize}
		\item $\P \equiv \coP$: si può sempre calcolare il complemento del risultato

		\item $\P \subseteq \NP \cap \coNP$, $\P \stackrel{?}{\not \equiv} \NP \cap \coNP$
	\end{itemize}
\end{itemize}
	\newpage
	
	\section{Algoritmi Probabilistici}
	% !TeX spellcheck = it_IT
% !TeX root = ../scompl.tex

\subsection{Algoritmi Montecarlo e Las Vegas}

Definizioni di:
\begin{itemize}
	\item Algoritmo probabilistico
	
	\item Algoritmi Montecarlo: tempo deterministico, risultato possibilmente sbagliato
	\begin{itemize}
		\item One-sided: sempre giusto quando esce 1 ($q(I) = A(I,Z) = 1$)
		
		\item Two-sided: può sbagliare su entrambi gli output
	\end{itemize}
	
	\item Algoritmi Las Vegas: risultato corretto, tempo non deterministico
\end{itemize}

Amplificazione:
\begin{itemize}
	\item Montecarlo one sided: si può ridurre la probabilità di errore con esecuzioni indipendenti dell'algoritmo
	
	\item Montecarlo two-sided: voto di maggioranza su esecuzioni indipendenti, usando Chernoff-Hoeffding il voto di maggioranza è "probabilmente" corretto
\end{itemize}

Da Las Vegas a Montecarlo One-sided: 
\begin{itemize}
	\item Trovare $t(n) \geq c \cdot f(n)$, con $f(n)$ tempo di esecuzione massimo atteso dell'algoritmo Las Vegas
	
	\item Blocca l'esecuzione se l'algoritmo non termina in $t(n)$ passi, restituendo 0
	
	\item Sempre corretto se torna 1, errore con probabilità $\leq \frac{1}{c}$ per disuguaglianza di Markov
\end{itemize}

\subsection{Classi di complessità probabilistiche}

$\BPP$: problemi risolti efficientemente da algoritmi Montecarlo two-sided
$$ \Pr \left(B (I,Z) \neq q(I)\right) \leq \frac{1}{3} = \frac{1}{2} - \frac{1}{p'(|I|)} $$
\begin{itemize}
	\item $\P \subseteq \BPP$, $\P \stackrel{?}{\equiv} \BPP$
	
	\item $\BPP \stackrel{?}{\equiv} \NP$
	
	\item $\BPP \equiv \coBPP$: la definizione è simmetrica rispetto al complemento
\end{itemize}

$\RP$: problemi risolti efficientemente da algoritmi Montecarlo one-sided
$$
\begin{array}{r l}
	q(I) = 1 \implies & \displaystyle \Pr \left(B(I, Z) = 1\right) \geq \frac{2}{3} = \frac{1}{p'(|I|)} \\
	q(I) = 0 \implies & \displaystyle \Pr \left(B(I, Z) = 0 \right) = 1 
\end{array}
$$
\begin{itemize}
	\item $\P \subseteq \RP$
	
	\item $\RP \subseteq \NP$: riscrivendo la def di $\NP$
	
	\item $\coRP \subseteq \coNP$, $\P \subseteq \coRP$
	
	\item $\RP \subseteq \BPP$, $\coRP \subseteq \BPP$
\end{itemize}

$\ZPP$: def come $\ZPP \equiv \RP \cap \coRP$, problemi risolti in tempo atteso limitato da un polinomio nella lunghezza dell'istanza da algoritmi Las Vegas
$$
\begin{array}{r l}
	q(I) = 1 \implies & \displaystyle \Pr \left(B(I, Z) = 1\right) \geq \frac{2}{3} \wedge \Pr (B' (I, Z') = 1) = 1 \\
	q(I) = 0 \implies & \displaystyle \Pr \left(B'(I, Z') = 0\right) \geq \frac{2}{3} \wedge \Pr (B' (I, Z') = 0) = 1
\end{array}
$$
\begin{itemize}
	\item Se $X \in \ZPP$, alg $A$ che su $I$ esegue $B$ e $B'$ terminando quando $B(I, Z) = 1$ o $B'(I, Z') = 0$, la prob che ciò non accada è $1/3$; usando il $\Ex$ della Geometrica, $3/2 < 2$ ripetizioni è il numero atteso
	
	\item $\ZPP \subseteq \RP$: si può passare da LV a MC-OS
	
	\item $\P \subseteq \ZPP$: dato che $\P \subseteq \RP$ e $\P \subseteq \coRP$
\end{itemize}

\subsection{Estrattore di Von Neumann}

Si vuole trasformare una sorgente non perfettamente casuale in una completamente casuale. Lancia una moneta truccata finché non si ottengono due valori diversi di seguito, le due sequenze possibili sono i due possibili risultati, hanno la stessa prob.

Numero di lanci atteso prima di una sequenza utile: 
\begin{itemize}
	\item Definiamo $Z_k = 1 \Leftrightarrow$ se i lanci $2k$ e $2k-1$ sono una sequenza utile
	
	\item Per valore atteso della Geometrica trovo il numero medio di lanci
\end{itemize}

\subsection{Il problema del Coupon Collector}

Vogliamo osservare gli $n$ possibili valori distinti, ognuno ha prob $\frac{1}{n}$. 

Vogliamo trovare il numero minimo di realizzazioni per osservare ciascun $a_i$ almeno una volta:
\begin{itemize}
	\item $N_1, \dots, N_n$, con $N_i$ numero di estrazioni aggiuntive per ottenere l'$i$-esimo valore distinto; sono tutte Geometriche
	
	\item Data la probabilità di osservare un nuovo valore calcolo 
	$$ \Ex [N] = \sum_{i = 1}^n N_i = \dots = n \left(\ln n + \Theta (1)\right) $$
\end{itemize}

\subsection{Reservoir Sampling}

Vogliamo una struttura dati che contenga $k$ elementi e dato uno stream in ingresso per ogni $t \geq k$, ognuno dei primi $t$ elementi è presente nella struttura con prob $\frac{k}{t}$. Modello streaming e spazio $\Theta (k)$.

Definizione dell'algoritmo.

Dimostriamo che $\forall t \geq k$ vale $\Pr\left(x_i \in R_t \right) = \frac{k}{t}$, $\forall i \leq t$:
\begin{itemize}
	\item Per induzione, base con $t = k$
	
	\item Step: con $t \geq k$ dimostriamo $\Pr \left(x_i \in R_{t+1}\right) = \frac{k}{t+1}$, $\forall i \leq t+1$
	\begin{align*}
		\Pr \left(x_i \in \R_{t+1}\right) & = \Pr \left(x_i \in R_t\right) \Pr (x_i \in R_{t+1} \mid x_i \in R_t ) \\ 
		& = \Pr \left(x_i \in R_t\right) \left( 1 - \Pr (x_i \notin R_{t+1} \mid x_i \in R_t) \right) \\
		& = \frac{k}{t} \cdot \left(1 - \frac{k}{t + 1} \frac{1}{k} \right) = \frac{k}{t+1}
	\end{align*}
\end{itemize}

\subsection{Algoritmo di Karger per il taglio minimo}

Dato un multigrafo non orientato, vogliamo il taglio che corrisponde a una partizione ammissibile con il costo minimo.

Definizione di Karger-Base e Karger.

Vogliamo prima trovare la prob che Karger-Base torni il taglio minimo: 
\begin{itemize}
	\item Prob che l'$i$-esimo arco non sia nel taglio $\Gamma^\ast$
	$$ p_i = \Pr \left(A_i \mid A_1, \dots A_{i-1}\right) = \dots \geq \frac{|V| - i - 1}{|V| - i + 1} $$
	
	\item Prob che l'alg restituisca $\Gamma^\ast$
	$$ \Pr (\text{ret } \Gamma^\ast) = \prod_{i = 1}^{|V| - 2} p_i \geq \dots = \frac{1}{\binom{|V|}{2}} $$
\end{itemize}

Possiamo amplificare l'alg per ottenere probabilità arbitrariamente bassa:
\begin{itemize}
	\item Probabilità che fallisca su $M$ esecuzioni indipendenti
	$$ \Pr (\text{fail}) = \Pr (\text{ret } \Gamma^\ast)^M \leq e^{- \frac{1}{\binom{|V|}{2}} M} = \epsilon$$
	scegliendo $M = \left\lceil \binom{|V|}{2} \ln \frac{1}{\epsilon} \right\rceil$
\end{itemize}

Il tempo totale di esecuzione è $\O \left(|E| |V|^2 \ln \frac{1}{\epsilon}\right)$.
	\newpage
	
	\section{Hashing e Randomizzazione}
	% !TeX spellcheck = it_IT
% !TeX root = ../scompl.tex

\subsection{Hashing Universale}

Sia $\U$ con $|\U| = m \gg 1$ l'insieme universo. Consideriamo una tabella $T$ di dim $1 < n \ll m$. Vogliamo memorizzare $S \subseteq \U$ nella tabella usando funzioni di hash. Ovviamente ci saranno collisioni, ma vogliamo minimizzarle
\[ \Pr \left(H(u) = H(v)\right) = \frac{1}{n}, \quad \forall u,v \in \U, \ u \neq v \tag{\dag} \]
\begin{itemize}
	\item Possiamo farlo tramite una famiglia $\H$ di funzioni $h : \U \rightarrow \left\{0, \dots, n-1 \right\}$ e usare una funz cas
	
	\item La prima condizione non basta, vogliamo che le funzioni occupino spazio al più $\Theta (\log m)$ e siano efficientemente calcolabili $(\ddag)$
\end{itemize}

Una famiglia che soddisfa entrambe le condizioni è una famiglia universale di funzioni di hash. Per dimostrarne l'esistenza:
\begin{itemize}
	\item Consideriamo $n = p$ con $p$ primo e $r = \lceil \log_2 m / \log_2 p \rceil$ il numero di cifre per rappresentare tutti i possibili valori
	
	\item Consideriamo $\H = \left\{h_a \mid a \in \A \right\}$, con $\A = \left\{0, \dots, p-1\right\}^r$ e $h_a : \A \rightarrow \left\{0, \dots, p-1\right\}$ def come $h_a (u) = \bm a \cdot \bm x$ con $\bm x = [u]_p$
\end{itemize}

Ogni $h_a$ è rappresentabile con $\lceil \log_2 |\A|\rceil = \Theta (r \log p) = \Theta (\log m)$ ed è calcolabile in modo efficiente (cond $(\ddag)$).

Per la condizione $(\dag)$:
\begin{itemize}
	\item Se $H$ è estratta a caso da $\H$ allora $\Pr \left(H(u) = H(v)\right) \leq \frac{1}{p}$
	
	\item Se $u \neq v$ vuol dire che almeno una cifra in base $p$ è diversa: supponendo siano diversi in posizione $j$ c'è solo un possibile valore per la cifra $j$ di $a \in \A$ che rende vera $h_a (u) = h_a (v)$
\end{itemize}

\subsection{Conteggio approssimato}

Vogliamo trovare in una tabella $A$ di $n$ interi gli elementi che si ripetono almeno $n/k$ volte ($n \gg k$). Risolviamo la versione approssimata, in cui ogni elemento che si ripete $n/k$ volte è nella lista finale, ma ogni elemento nella lista compare almeno $n/k - \epsilon n$ volte in $A$.

Definizione di Count-Min Sketch, con $\ell$ tabelle di hash di dimensione $b$.

Probabilità di errore della struttura con funzioni estratte casualmente da una famiglia universale $\H$:
\begin{itemize}
	\item Troviamo il valore atteso per il contatore di ogni $x$ per ogni funzione di hash $h_i$:
	$$\Ex [Z_i] = \Ex \left[n_x + \sum_{y \neq x} n_y \Ind \left\{H_i (y) = H_i (x)\right\}\right] = \dots \leq n_x + \frac{n}{b} $$
	
	\item Troviamo la probabilità che il contatore di una funzione "sfori" di più di $\epsilon n$:
	$$ p_i = \Pr \left(Z_i \geq n_x + \epsilon n\right) = \dots \leq \frac{1}{e} $$
	
	\item Troviamo la probabilità che tutte le funzioni di hash "sforino". 
	$$ c_i = \Pr \left(\text{\textsc{Count}}(X) \geq n_x + \epsilon n \right) = \prod_{i = 1}^\ell p_i \leq e^{-\ell} $$
	
	\item Troviamo la probabilità che un qualsiasi elemento della tabella abbia conteggio "sforato"
	$$ \Pr \left(\exists x \in A: \text{\textsc{Count}}(x) \geq n_x + \epsilon n\right) \leq \sum_{x \in A} c_i \leq n e^{-\ell} \leq \delta $$
\end{itemize}
Fissando $\delta = 1/100$ abbiamo confidenza del 99\% e $b = \left(\frac{1}{\epsilon}\right)$, $\ell = \Theta (\log n)$.

\subsection{Proiezioni casuali}

Vogliamo stimare la distanza Euclidea in $\R^d$ limitando la complessità computazionale, ad esempio per calcolare Nearest Neghbor dei punti in $S \subseteq \R^d$. 

Mostriamo che per ogni $0 < \epsilon$, $\delta < 1$ esiste $k = \O \left(\frac{1}{\epsilon^2} \ln \frac{|S|}{\delta} \right)$ e una classe $\F$ di funzioni $f: \R^d \rightarrow \R^k$ tale che
$$ \left(1 - \epsilon\right) \|\bm x - \bm x' \|^2 \leq \| f(\bm x) - f (\bm x') \|^2 \leq (1 + \epsilon) \| \bm x - \bm x' \|^2, \ \bm x, \bm x' \in S $$
con prob almeno $1 - \delta$ rispetto all'estrazione di $f \in \F$.

Usiamo $k$ funzioni casuali rappresentate da $k$ vettori casuali $\bm Z_1, \dots, \bm Z_k \in \R^d$ t.c. $ f_j (\bm x) = \bm Z_j^\top \bm x$. I vettori $\bm Z_j$ sono ottenuti generando ciascuna componente $Z_{j,i}$ con estrazioni indipendenti da una distr di prob con media 0 e varianza 1.

Verifichiamo che ogni $f$ si possa usare per stimare la distanza Euclidea tra due punti: 
$$ \Ex \left[\left(f_j (\bm x) - f_j (\bm x')\right)^2 \right] = \dots = \|\bm x - \bm x' \|^2 $$

Definiamo la proiezione casuale $f: \R^d \rightarrow \R^k$: 
$$ f(\bm x) = \left(\frac{f_1 (\bm x)}{\sqrt{k}}, \dots, \frac{f_k (\bm x)}{\sqrt{k}}\right) $$
e da notare che è una trasformazione lineare, quindi anch'essa stima la distanza quadrata, in quanto si tratta della media di $k$ stimatori di tale distanza.

Def $\bm v = \bm x - \bm x'$. Assumendo che le $Z_{i,j}$ abbiano una distr Normale, per ogni $\epsilon, \delta > 0$ e per ogni $\bm v$ t.c. $\|v\| = 1$
$$ \Pr \left(\left|\|f (\bm v)\|^2 - 1\right| > \epsilon \right) \leq \delta $$
e sappiamo che 
$$ \|f(\bm v)\|^2 = \frac{1}{k} \sum_{j = 1}^k \left(\bm Z_j^\top \bm v \right)^2 $$
Notando che le var cas $V_j = \left(\bm Z_j^\top \bm v\right)^2$ sono i.i.d. con media $\mu = 1$ e possiamo riscrivere la prob sopra 
$$ \Pr \left(\left|\frac{1}{k} \sum_{j \in [k]} V_j - \mu \right| > \epsilon \right) \leq e^{- \O (k \epsilon^2)} $$
Consideriamo gli eventi 
$$ A_{\bm x, \bm x'} = \left|\frac{\|f(\bm x) - f (\bm x') \|^2}{\| \bm x - \bm x' \|^2} - 1\right| > \epsilon $$
Troviamo la prob 
$$ \Pr \left(\exists \bm x, \bm x' \in S : A_{\bm x, \bm x'}\right) = \dots \leq n^2 \delta $$
per $k = \O \left(\frac{1}{\epsilon^2} \ln \frac{1}{\delta} \right)$.

Da questo ne deduciamo che vale 
$$ (1 - \epsilon) \| \bm x - \bm x' \|^2 \leq \| f(\bm x) - f(\bm x') \|^2 \leq (1 + \epsilon) \| \bm x - \bm x' \|^2, \quad \forall \bm x, \bm x' \in S $$
con prob almeno $1-\delta$ rispetto all'estrazione delle $Z_{j,i}$.
	\newpage
	
	\section{Clustering e Randomizzazione}
	% !TeX spellcheck = it_IT
% !TeX root = ../scompl.tex

\subsection{Correlation Clustering}

Caso speciale del clustering. Si ha una funzione di similarità $\sigma: \E \rightarrow \{-1, +1\}$, con $V \equiv [n]$ nodi in input, $\E \equiv \binom{n}{2}$ insieme di coppie di punti distinti. Un clustering $\C$ è una partizione di $V$ in cluster $C_i$, $i \in [k]$. Dati $\C$ e $\sigma$, $\Gamma_\C$ contiene tutte le coppie che causano errore, ovvero t.c. $\sigma (u,v) = -1$ con $u,v \in C_i$ o t.c. $\sigma(u,v) = +1$ con $u \in C_i$ e $v \in C_j$ ($i \neq j$).

Un triangolo è qualunque tripla non ordinata $T = \left\{u,v,w\right\}$. Un bad triangle BT è un triangolo con segni dei valori sui lati del triangolo $\left\{+,+,-\right\}$. $\T$ è l'insieme di tutti i BT di $V$.

Definizione di KwikCluster.

\paragraph{Lemma:} La somma pesata di tutti i BT è lower bound per OPT: se la somma dei pesi $\left\{\beta_T \geq 0 \mid T \in \T \right\}$ incidenti su ogni lato è t.c. $\sum_{T \in \T(e)} \beta_T \leq 1$, $\forall e \in \E$, allora $\sum_{T \in \T} \beta_T \leq \Opt$.

Cerchiamo un bound per l'errore atteso di KwikCluster. Sia $\Gamma_A$ l'insieme di lati sbagliati per clustering emesso e $V_r$ l'insieme dei nodi alla chiamata ricorsiva $r$.

\paragraph{Lemma:} un lato $e \in \E$ causa errore $e \in \Gamma_A$ sse esiste chiamata $r$ e triangolo $T \in \T$ t.c. $T \subset V_r$, $T \in \T(e)$ e il pivot del round fa parte di $T$, $\pi_r = T \setminus e$. Per dimostrarlo: 
\begin{itemize}
	\item $e \in \Gamma_A \implies$ esistono $r$, $T \subseteq V_r$ t.c. $\pi_r = T \setminus e$ e $T \in \T(e)$
	
	\item Per l'inverso, assumiamo $e = \left\{u,v,\right\}$, $T = \left\{u, \pi_r, w\right\} \subseteq V_r$. Analizzare entrambi i possibili casi
\end{itemize}
Quindi all'iterazione $r$ viene fatto errore esattamente su uno dei lati di ogni $T \in \T$. Chiamando $A_T$ l'evento $\left\{ \exists r \mid T \subseteq V_r \wedge \pi_r \in T \right\}$. Per $T,T' \in \T(e)$, $A_T$ e $A_{T'}$ non possono accadere assieme, e il valore atteso dato dalla sequenza casuale di pivot 
$$ 1 = \sum_{T \in \T(e)} \Ind \left\{A_T \wedge e \in \Gamma_A \right\} = \sum_{T \in \T(e)} \Pr \left(e \in \Gamma_A \mid A_T \right) \Pr \left(A_T\right) = \sum_{T \in \T(e)} \frac{1}{3} \Pr (A_T) $$
Applicando il primo lemma con $\beta_T = \frac{1}{3} \Pr (A_T)$ per ogni $T \in \T$ si ottiene che $\Ex\left[|\Gamma_A|\right] \leq 3 \Opt$.

\subsection{$k$-Means}

Vogliamo partizionare l'insieme finito $\X \subset \R^d$ di punti in $k>1$ cluster. Il costo di un punto $\bm x \in \X$ all'interno di una partizione $\C$ è la sua distanza Euclidea dal centro del cluster più vicino $\phi(\C, \bm x)$. Il costo della partizione $\Phi(\C)$ è la somma del costo di tutti i punti. Il problema assume implicitamente che i punti in $\X$ siano campionati da $k$ distribuzioni Gaussiane sferiche, le quali medie sono i centri e le quali varianze upper bound per il costo ottimale.

Algoritmo di Lloyd. Il tempo per iterazione è $\O (nkd)$, si può ridurre a $\O(n k \ln n)$ con proiezioni casuali, ciò moltiplica $\Opt$ per una costante. Il numero di iterazioni worst case è $2^{\Omega\left(\sqrt{n}\right)}$.

Vogliamo dimostrare che non approssima $\Opt$ per nessuna costante:
\begin{itemize}
	\item Per ogni costante, def istanze 1-dimensionali del problema con $k = 3$ t.c. $|\X| = n$ e $n-2$ punti sono in $[0,1]$ equispaziati, gli altri 2 sono a distanza $2 \sqrt{an}$ e $3 \sqrt{an}$
	
	\item La probabilità che l'alg non scelga entrambi gli outlier tende a $p_n \rightarrow 1$ per $n \rightarrow \infty$
	
	\item Nel caso in cui non vengano scelti entrambi gli outlier $\Phi (\C)/\Opt = \Omega (a)$
\end{itemize}

Vogliamo mostrare che, se i centri si muovono, il valore di $\Phi$ decresce strettamente:
\begin{itemize}
	\item Definiamo la funzione $\psi$ per il costo della partizione dati Cluster e centri (non per forza corrispondenti)
	
	\item L'alg assegna ogni punto al cluster più vicino, se $\bm c_i \neq \bm c_i$ per qualche $i$, allora $\psi$ decresce strettamente (dopo l'aggiornamento)
\end{itemize}

Possiamo anche vedere che l'algoritmo termina in al più $k^{|\X|}$ iterazioni per qualsiasi input, in quanto la funzione $\Phi$ è strettamente decrescente e può assumere al più $k^n$ valori distinti.

\subsection{$k$-Means$++$}

Definizione dell'algoritmo.

\paragraph{Lemma:} Per ogni $A \in \C^\ast$ e per ogni $i \in [k]$: $\Ex \left[\phi (\C_i, A) \mid \bm c_i \in A, \C_{i-1} \right] \leq 8 \phi (\C^\ast, A)$. Per dimostrarlo:
\begin{itemize}
	\item Troviamo il valore atteso del costo della partizione dopo aver scelto il primo centro $\Ex \left[\phi \left(\C_1, A\right) \bm c_1 \in A\right]$, ovvero la media del costo di ogni punto (scelto u.a.r.), in totale $\leq 2 \phi (C^\ast, A)$
	
	\item Sapendo che $\phi (\C_{i-1}, \bm a) \leq 2 \left(\phi (\C_{i-1}, \bm x) + \| \bm a - \bm x \|^2 \right) $, possiamo farne la media su tutti i punto $\bm x \in A$
	
	\item Inoltre per ogni $\bm x \in \X$ vale che $\phi (\C_i, \bm x) = \min \left\{\phi (\C_{i-1}, \bm x), \|\bm x - c_i\|^2 \right\}$
	
	\item Calcoliamo il valore atteso all'aggiunta di un nuovo centro, ovvero il valore cercato nella tesi. Tale valore atteso è la prob che venga scelto $\bm a \in A$ per il costo di $\bm a$, per ogni $\bm a$, i due termini si possono dividere, sostituire come indicato, fino ad arrivare a 4 volte il valore atteso trovato all'inizio della dimostrazione
\end{itemize}

Facciamo delle assunzioni: per ogni $A \in \C^\ast$
\begin{itemize}
	\item $\phi(\C^\ast, A) = 1$
	
	\item Per ogni $i \in [k]$, se $A$ è coperto in $\C_i$ allora $\phi(\C_i, A) = 1$, altrimenti $\phi(\C_i, A) = L$
\end{itemize}

\paragraph{Lemma:} Considerando le assunzioni $\Ex\left[\Phi(\C)\right] \leq \left(2 + \ln k\right) \Opt $. Per dimostrarlo:
\begin{itemize}
	\item Osserviamo che $\Phi (\C_k) = \Phi (\C_0) + \sum_{i = 0}^{k-1} \left(\phi(\C_{i+1} - \Phi (\C_i))\right)$, ovvero il costo finale è pari a costo iniziale sommato alla variazione a ogni passo, la quale può essere $(L-1)$ se viene coperto un cluster $A \in \C^\ast$ precedentemente scoperto, $0$ altrimenti. Consideriamo il valore atteso di tale valore
	
	\item Troviamo la probabilità che all'iterazione $i+1$ venga scelto un centro in un qualsiasi $A$ scoperto: $p_{i+1} = \frac{N_i L}{N_i L + (k- N_i)} \geq \frac{(k-i) L}{(k-i)L + i}$
	
	\item Data la prob, possiamo trovare la variazione di costo attesa per ogni iterazione (la parte dentro la sommatoria): $< \frac{k}{k-i}$
	
	\item Possiamo calcolare il costo totale, ricordando il bound sulla somma armonica e che sotto le assunzioni $\Opt = \Phi (\C^\ast) = k$
\end{itemize}
	\newpage
	
	\section{Giochi e Mercati}
	% !TeX spellcheck = it_IT
% !TeX root = ../scompl.tex

\subsection{Hedge e Exp3 per decision-making sequenziale}
	
	
	
\end{document}
